\section{Inverse basins in every arithmetic progression}\label{sec:inverse}

This section fixes the forcing $d=(1,0,-1)$, with phase zero immediately
before $+1$. Write $F(n)=F_{d,0}^{(3)}(n)$ for its return map. A target
$(v,i)$ is safe precisely when $v\ge0$ for $i=0$ and $v\ge1$ for $i=1,2$.
All predecessors below refer to this fixed forcing, with phase included.

\begin{theorem}[Inverse growth and arithmetic progressions]\label{thm:inverse-growth}
There is a number $\delta$, with $0.535<\delta<0.536$, such that the following
holds. For every safe target $(v,i)$, every modulus $M\ge1$, every residue
$a\bmod M$, and every launch phase $j$, there are constants $c>0$ and $X_0$
such that, for $X\ge X_0$, at least $cX^\delta$ positive integers $x\le X$
satisfy $x\equiv a\pmod M$ and their trajectories from phase $j$ reach
$(v,i)$. The constants may depend on the target, progression, and phase.
The exponent is defined exactly by
\begin{equation}\label{eq:inverse-exponent}
 \sum_{o=0}^{8}m_o(2^{24}/3^o)^{-\delta}=1,
 \qquad (m_0,\ldots,m_8)=(1,4,17,35,57,69,53,32,13).
\end{equation}
\end{theorem}

An exact finite certificate supplies the branches of a growing inverse
tree. A bounded-length selector sends its vertices into the required
progression with bounded size cost and multiplicity.

For ordinary Collatz, Monks proved that every positive orbit merges
with an orbit starting in any prescribed progression, possibly after
different numbers of steps \cite[Theorem~1.1]{monks-2006}.
Theorem~\ref{thm:inverse-growth} counts predecessors of a specified
target at a specified phase for the forced map.

\subsection{Elementary inverse branches and eligible targets}

The full three-step return has eight affine pieces.  In the table below,
the parity string is read in chronological order from the three inputs
to the shortcut map; it is not the binary expansion of the residue $r$.
Every row is valid for all integers $t\ge0$.
\begin{center}
\begin{tabular}{rrl|rrl}\toprule
$r$&parities&$F(8t+r)$&$r$&parities&$F(8t+r)$\\\midrule
0&010&$3t$&4&011&$9t+7$\\
1&111&$27t+7$&5&110&$9t+6$\\
2&001&$3t+1$&6&000&$t$\\
3&101&$9t+4$&7&100&$3t+2$\\\bottomrule
\end{tabular}
\end{center}
For example, the residue $r=3$ follows
\[
 8t+3\ \longmapsto\ 12t+6\ \longmapsto\ 6t+3
 \ \longmapsto\ 9t+4.
\]
The impulses are applied after the shortcut update at each arrow.
The slopes in the table range from $1$ to $27$ as functions of $t$.

Four convenient rows of the complete table give
\begin{equation}\label{eq:inverse-four}
 F(8t+6)=t,\quad F(8t)=3t,\quad
 F(8t+2)=3t+1,\quad F(8t+7)=3t+2.
\end{equation}
Consequently the following are inverse branches:
\[
 A(y)=8y+6,\quad B_0(y)=\frac{8y}{3},\quad
 B_1(y)=\frac{8y-2}{3},\quad B_2(y)=\frac{8y+5}{3}.
\]
Here $A$ is always applicable, and $B_r$ is applicable precisely when
$y\equiv r\pmod3$. For $y\ge2$ every applicable branch strictly increases
$y$ and its output $x$ satisfies $x+1\le8(y+1)$. Call a phase-zero value
\emph{eligible} if $y\ge2$ and $3\nmid y$.

Every phase-zero $u\ge0$ has an eligible predecessor two laps earlier.
Explicitly, choose
\begin{equation}\label{eq:eligible-attachment}
 y=\begin{cases}
 64u+54,&3\nmid u,\\
 64t+6,&u=3t,\ 3\nmid t,\\
 64t+16,&u=3t,\ 3\mid t.
 \end{cases}
\end{equation}
The formulas in \eqref{eq:inverse-four} prove $F^2(y)=u$ in each case;
$y\ge2$ and $3\nmid y$ are immediate. For a target $(v,1)$ or $(v,2)$,
the phase-zero values $2(v-1)$ and $4v-2$, respectively, reach that target
after one or two updates. These values are nonnegative on the safe domain.
Thus \eqref{eq:eligible-attachment} attaches every safe augmented target.

For instance, phase-zero targets $0,1,3,9$ have the eligible two-lap
predecessors $16,118,70,208$, respectively. Their exact returns are
\[
 16\mapsto6\mapsto0,\quad118\mapsto14\mapsto1,\quad
 70\mapsto8\mapsto3,\quad208\mapsto78\mapsto9.
\]
Each arrow in this display is one whole three-update lap. The third
case of \eqref{eq:eligible-attachment} matters even for target zero:
simply using $A^2(0)=54$ would give a parent divisible by three,
outside the eligible class used for the growing tree.

In the ordinary inverse graph a positive multiple of three has only
its iterated doublings as predecessors. The backtracing theorem of
Monks et al. assumes a target coprime to three
\cite[Theorem~4.1]{monks-et-al-2013}. Here
\eqref{eq:eligible-attachment} and the phase lifts cover every safe
target, including zero.

\subsection{The finite 24-step certificate}

Applegate and Lagarias enumerate ordinary inverse trees by depth,
residue class and odd-step count \cite[Sections~2--3]{applegate-lagarias-1995-i}.
Krasikov and Lagarias refine the estimates by residue-dependent
difference inequalities \cite[Theorems~2.2 and~6.1]{krasikov-lagarias}.
Here we enumerate forced branches, bound their sizes, and prove that
the branches counted in the inverse tree give distinct predecessors.

Put $S=2^{24}$. For an itinerary of 24 input parities with $o\le8$ odd
inputs, Proposition~\ref{prop:affine-parity} gives a unique initial residue
$r\bmod S$. Its inverse formula is
\begin{equation}\label{eq:inverse-offset}
 x=\frac{Sy+c}{3^o},\qquad c=3^or-SF^8(r).
\end{equation}
The row with no odd inputs is especially transparent: apply $A$ eight
times to obtain
\[
 A^8(y)=8^8y+6(1+8+\cdots+8^7)
       =2^{24}y+14{,}380{,}470.
\]
This verifies the first offset without an enumeration and explains why
24 literal updates correspond to eight returns of $F$.
Integrality restricts $y$ to one class modulo $3^o$. To count children
coprime to three, inspect its three lifts modulo $3^{o+1}$, retain those
with $3\nmid x$, and discard parents divisible by three. Distinct formulas
have distinct residues $r$, so their integer children are distinct.

\begin{lemma}[Certified branch minima]\label{lem:inverse-certificate}
For each $y\ge34$ with $3\nmid y$, and each $0\le o\le8$, there are at
least $m_o$ distinct children $x$ such that
\[
 F^8(x)=y,\qquad 3\nmid x,\qquad x>y,\qquad
 x+1\le(S/3^o)(y+1).
\]
\end{lemma}
\begin{proof}[Exact finite verification]
The complete enumeration and its output are
\path{certificates/inverse/certify_growth.py} and
\path{certificates/inverse/growth_certificate.json}. The algorithm enumerates
all $\sum_{o=0}^8\binom{24}{o}=1{,}271{,}626$ parity words with at most
eight odd inputs by composing inverse updates. For every formula it
recovers $r=c(3^o)^{-1}\bmod S$, rejects repeated $r$, independently replays
all 24 forward updates at $r$, and checks the odd count and offset.
It then counts the admissible children on every unit parent class modulo
$3^{o+1}$. The minima and offset extremes are as follows.
\begin{center}\small
\begin{tabular}{rrrrr}\toprule
$o$ & $m_o$ & Parity words & Minimum $c$ & Maximum $c$\\\midrule
0&1&1&14380470&14380470\\
1&4&24&1198370&15878434\\
2&17&276&-25765018&16440166\\
3&35&2024&-40594894&16650802\\
4&57&10626&-70928746&16729750\\
5&69&42504&-147774526&16759234\\
6&53&134596&-197825722&16769926\\
7&32&346104&-300203566&16772842\\
8&13&735471&-559561354&16770655\\\bottomrule
\end{tabular}
\end{center}
For each enumerated offset the checker verifies, using integers,
\[
 34(S-3^o)+c>0,\qquad c+3^o\le S.
\]
These inequalities imply strict growth and the size bound for every
$y\ge34$. A positive integral $x$ with residue $r\in[0,S)$ is $r+Sk$
with $k\ge0$, so the forward affine identity applies. Admissibility of the
forcing guarantees that its intermediate states are safe. The enumeration
and residue grids are finite and exhaust the indicated cases. This is a
computer-assisted finite lemma, not a Lean theorem in this release; the
hash in the output is an integrity aid, not the proof of its arithmetic.
\end{proof}

Set $L_o=S/3^o>1$. The left side of \eqref{eq:inverse-exponent} decreases
continuously from 281 to zero, establishing a unique positive root.
At exponent $1/2$ its value is
\[
 \frac{3049+1594\sqrt3}{4096}>\frac{4643}{4096}>1,
\]
so $\delta>1/2$ without numerical approximation. The reported interval
uses integer powers: with $G=10^9$, find $a_o$ such that
\[
 a_o^q S^p\le G^q3^{op}<(a_o+1)^q S^p.
\]
Then $a_o/G$ and $(a_o+1)/G$ bound $L_o^{-p/q}$ from below and strictly
from above. Weighted sums certify mass greater than one at $p/q=107/200$
and less than one at $67/125$. The value
$0.535422959115049\ldots$ is an approximation, not the definition.

Keeping the different size factors is essential to this exponent. If all
281 selected children were charged the common worst factor $S$, the
resulting exponent would be $\log(281)/\log(S)<1/2$, since
$281<\sqrt S=4096$. The certified gains come from the many children
with factors $S/3^o$ substantially smaller than $S$. Formula
\eqref{eq:inverse-exponent} retains those gains while still giving a
uniform argument at every eligible parent; it is not an estimate from
an average sampled trajectory.

\begin{lemma}[Weighted inverse counting]\label{lem:inverse-count}
For every eligible root $y\ge34$, its phase-zero backward orbit contains
at least $((X+1)/(S(y+1)))^\delta$ distinct values at most $X$, whenever
$X+1\ge S(y+1)$.
\end{lemma}
\begin{proof}
At every vertex retain exactly $m_o$ children of each type $o$, assigning
each edge weight $L_o^{-\delta}$. The outgoing weights sum to one.
Different paths of equal depth cannot collide, because $F^8$ determines
the parent uniquely. A collision of depths $a<b$ would, on applying $a$
forward blocks, identify the root with its own descendant of depth $b-a$.
Strict increase along inverse edges rules this out as well.

Write $W=(X+1)/(S(y+1))\ge1$. Stop a path when its product of edge factors
first reaches $W$. Since the smallest factor exceeds one, the stopping
tree has finite depth. Its stopped products lie between $W$ and $SW$;
the corresponding integers satisfy $x+1\le SW(y+1)=X+1$.
The weights of all stopped leaves sum to one, while each is at most
$W^{-\delta}$. Hence there are at least $W^\delta$ distinct leaves.
For $W=1$, stopping at the root gives the same assertion.
\end{proof}

\subsection{Ternary steering with uniform finite choices}

The modular translation and admissible backtracing strategy has direct
antecedents in Monks and Monks et al.
\cite{monks-2006,monks-et-al-2013}. The latter's Theorem~4.1 bounds the
number of \emph{odd inverse operations} needed to reach a progression,
independently of the target. Its convention $l(s)$ counts those operations,
whereas $\|s\|$ also counts the intervening even operations.
For the counting transfer here we need a finite selector with bounded
\emph{total physical length}, size cost, and multiplicity. The following
proofs establish these properties directly, including the integral domains
of the modular words; a bound on odd inverse count alone would not suffice.

\begin{lemma}[Ternary steering]\label{lem:ternary-steering}
For each $k\ge0$, every eligible target $q$ has an inverse predecessor
in every class modulo $3^k$, using a uniformly bounded number of branches
from $\{A,B_0,B_1,B_2\}$. The bound depends only on $k$.
There are also arbitrarily large such predecessors, if the length is
allowed to increase.
\end{lemma}
\begin{proof}
For $k=0$, the identity is a bounded-length choice, and the iterates
$A^h(q)$ supply arbitrarily large choices. Hence assume $k\ge1$.
The identities
\[
 A^{2h}(q)=\frac{64^h(7q+6)-6}{7},\qquad
 v_3(64^h-1)=2+v_3(h)\quad(h\ge1)
\]
show that the $A^2$ orbit of a unit $q$ modulo $3^k$ consists of all lifts
of its residue modulo 9, with period $3^{k-2}$ for $k\ge2$.
For completeness the valuation identity follows by writing $64=1+9\cdot7$:
if $3\nmid h$, the first binomial term has valuation two and the later
terms have larger valuation; replacing an exponent $h$ by $3h$ increases
the valuation by one, since $a^2+a+1$ has valuation one when $a\equiv1\pmod9$.
The differences of two orbit points therefore have valuation exactly
$2+v_3(h_1-h_2)$. There are $3^{k-2}$ distinct points, exactly the number
of lifts of one residue modulo 9.

If $q\equiv5$ or $7\pmod9$, one application of $A$ moves it to residue
$1$ or $8$, respectively. On residues $1,2,4,8\pmod9$, respectively use
$B_1,B_2,B_1,B_2$; their outputs are units congruent to $2,1,1,2\pmod3$.
Varying the input along its $A^2$ orbit modulo $3^{k+1}$ and then applying
that $B$ branch supplies all lifts of one unit class modulo three at
precision $3^k$. An optional $A$, which is invertible and swaps the two
unit classes modulo three, supplies the other class. This proves steering
to any unit residue at any precision, with a finite bound for fixed $k$.

For any residue $a\bmod3^k$, first steer to the unit
$z\equiv(3a+2)8^{-1}\pmod{3^{k+1}}$ and apply $B_1$.
Its output is $a\bmod3^k$. All inputs can be kept at least two; applicable
branches then increase them. To produce this $z$ at precision $3^{k+1}$,
use the $A^2$ orbit at precision $3^{k+2}$; its period is $3^k$.
An exponent $0\le h\le3^k-1$ therefore suffices. The initial adjustment
by $A$, the selected $B$ branch, the optional $A$, and the final $B_1$
cost at most four additional branches. Thus the total length is at most
$2(3^k-1)+4=2\cdot3^k+2$. Taking additional full periods instead
gives arbitrarily large outputs. This latter assertion does not have a
bounded length requirement.
\end{proof}

\begin{lemma}[Exact domains of positive inverse words]\label{lem:inverse-domains}
A word $w$ of $\ell$ branches from $\{B_0,B_1,B_2\}$ has affine form
$B_w(y)=(8^\ell y+C_w)/3^\ell$. There is exactly one input residue modulo
$3^\ell$ on which every intermediate branch is integral and applicable.
Given any input residue $a\bmod3^k$, a word of $k$ letters can be chosen
whose domain is exactly that residue.
\end{lemma}
\begin{proof}
Composition proves the affine form. Final integrality gives a unique
residue, since 8 is invertible modulo $3^\ell$. Reducing that congruence
modulo $3^t$ gives the integrality condition for the first $t$ operations:
in the remaining affine composition its prefix numerator is multiplied
by a power of 8 and all later contributions have a factor $3^t$.
Thus all prefixes are integral. At each step exactly one of the three
formulas has integral value on a given residue modulo three; integrality
therefore is the stated branch condition. To get a word for $a$, start
with a sufficiently large positive representative of $a$ and choose its
applicable branch successively $k$ times. The unique domain contains $a$.
\end{proof}

For a concrete instance of the domain lemma, apply $B_2$ and then $B_1$.
The composed formula is
\[
 B_1(B_2(y))=\frac{64y+34}{9},\qquad y\equiv2\pmod9.
\]
Targets $y=2$ and $y=11$ give the legal inverse histories
$2\mapsto7\mapsto18$ and $11\mapsto31\mapsto82$.
Indeed $F^2(18)=2$ and $F^2(82)=11$. At $y=5$ the first branch gives
$B_2(5)=15$. The second requires residue $1\pmod3$, and the composed fraction
is nonintegral.

\begin{lemma}[Simultaneous odd-modulus steering]\label{lem:odd-steering}
Let $k\ge0$, $m\ge1$, and $\gcd(m,6)=1$. Every eligible target has a
predecessor in any specified class modulo $3^km$. For fixed $k,m$, the
choice can be made from a finite list of inverse words of bounded length.
Arbitrarily large choices also exist when length is unrestricted.
\end{lemma}
\begin{proof}
For $m=1$ use Lemma~\ref{lem:ternary-steering}. Otherwise all $B_r$ act as
permutations modulo $m$ with common linear coefficient $8/3$. Let $e$ be
the multiplicative order of $8/3$ modulo $m$, and put $Q=me$. Raising a
$B_r$ to power $e$ gives a translation, and raising a translation to power
$m$ gives the identity. Hence $B_r^Q$ is the identity modulo $m$, and its
inverse is represented by the positive word $B_r^{Q-1}$.

Modulo $m$ one has $B_1^{-1}B_0(y)=y+1/4$. Four repetitions give
translation by one. Replacing every inverse letter by $Q-1$ positive
copies produces a positive word $P$ of length $4Q$ representing $y+1$.
Let $w_0$ be the $k$-letter word from Lemma~\ref{lem:inverse-domains}
whose domain is the prescribed class $a\bmod3^k$. In application order
concatenate $w_0$, a positive representative of $w_0^{-1}$ modulo $m$,
and $P$, obtaining a word $V$ of length $\ell=Q(k+4)$.
It acts as translation by one modulo $m$, and every admissible use of it
starts in class $a\bmod3^k$.

The word $V^m$ has a unique admissibility class modulo $3^{m\ell}$.
Use ternary steering to obtain a predecessor $z\ge2$ of the original
target in that class. All letters of $V^m$ are then applicable.
The block boundaries $z,V(z),\ldots,V^{m-1}(z)$ start admissible copies
of $V$, so they all lie in class $a\bmod3^k$ and run through every
residue modulo $m$. Select the required one. For fixed moduli, the finite
ternary choice and at most $m-1$ copies of $V$ give a uniform length bound.
Using arbitrarily large ternary choices proves the separate unbounded-size
assertion. This domain argument is necessary: the modular permutation
calculation alone does not license arbitrary integer applications of $B_r$.
\end{proof}

\subsection{Binary digits, phases, and preservation of the exponent}

\begin{lemma}[Finite selector]\label{lem:finite-selector}
Fix a progression $a\bmod M$ and a launch phase $j$. There are constants
$D,C$ and finitely many rational affine formulas such that every eligible
phase-zero $y$ has an integer predecessor $x>0$ in that progression, at
phase $j$, given by one of these formulas on its exact admissibility
domain, reaching $y$ within $D$ physical updates and satisfying
$x+1\le C(y+1)$.
\end{lemma}
\begin{proof}
Write $M=2^\alpha3^bm$ with $\gcd(m,6)=1$. Choose a positive
representative $x_0$ of the progression. Choose $K\ge\alpha$ such that
launch phase $j$ reaches phase
zero in $K$ updates. Let its endpoint be $A_0$ and let $o$ count the odd
inputs in these $K$ updates. By the finite parity formula, for $t\ge0$,
\begin{equation}\label{eq:prefix-lift}
 x_0+2^K3^bmt\ \longmapsto\ A_0+3^{o+b}mt.
\end{equation}
The difference at each odd step is multiplied by $3/2$, at each even
step by $1/2$, and all impulses cancel from the difference. All sources
in \eqref{eq:prefix-lift} lie in the requested progression.

Choose a fixed $R$ such that $A^R(2)\ge\max(A_0,2)$, and first replace
$y$ by $A^R(y)$, still an eligible predecessor of $y$.
Lemma~\ref{lem:odd-steering} supplies a predecessor $z$ in class
$A_0\bmod3^{o+b}m$ with uniformly bounded inverse length. Branches grow
on inputs at least two, so $z\ge A_0$. Solving \eqref{eq:prefix-lift}
for $t$ gives a nonnegative integer and a positive $x$ in the progression.
The finite alphabet and length bound give finitely many affine formulas.
Their positive slopes and constant offsets give a common size constant
$C$ on $y\ge2$. The total physical length, including the fixed prefix
and padding, has a common bound $D$.
\end{proof}

\begin{proof}[Proof of Theorem~\ref{thm:inverse-growth}]
Attach an eligible phase-zero predecessor to the target using
\eqref{eq:eligible-attachment} and the phase lifts. If necessary, two
applications of $A$ put it above 34, since $A^2(2)=182$. Apply
Lemma~\ref{lem:inverse-count} to this root. Its at least $c_0Y^\delta$
eligible tree vertices below $Y$ each admit a selector predecessor
$x\le C(Y+1)-1$. A fixed output $x$ can be assigned to at most $D+1$
different tree vertices: they must appear among its first $D+1$ augmented
forward states. Thus there are at least $c_0Y^\delta/(D+1)$ distinct
outputs. Replacing $Y$ by $(X+1)/C-1$ proves the required lower bound
for all sufficiently large $X$.
\end{proof}

\begin{corollary}[A conditional route from exceptional-set bounds]
\label{cor:exceptional-set}
For this three-step forcing, fix a finite catalog of safe cycles and let
$E_j$ be the positive seeds at phase $j$ whose orbits never enter the
catalog. If one such exceptional safe trajectory exists, every progression
at every launch phase contains at least $cX^\delta$ exceptional seeds
below $X$ for all sufficiently large $X$. Consequently, a bound
$\#(E_j\cap\{x\le X:x\equiv a\pmod M\})=o(X^\delta)$ in even one
whole progression implies global capture. A bound $O(X^{1/2})$ suffices.
\end{corollary}
\begin{proof}
Every predecessor of an exceptional target is exceptional, since entry
into a cycle is forward invariant. Apply Theorem~\ref{thm:inverse-growth}
to that target. The claimed upper bound would contradict the resulting
lower bound. Since $\delta>1/2$, $O(X^{1/2})=o(X^\delta)$.
\end{proof}

The required exceptional-set upper bound remains open. Disjoint basins
can both have the stated inverse growth. Krasikov and Lagarias obtain
larger lower exponents for ordinary Collatz \cite{krasikov-lagarias}.

Tao's ordinary-map theorem concerns minima below any prescribed threshold
tending to infinity, for a set of logarithmic density one
\cite[Theorem~1.3 and Remark~1.4]{tao}. It supplies no exceptional-set upper
bound of the form required by Corollary~\ref{cor:exceptional-set} for this
forced map. Even natural-density-one capture would give only $o(X)$
exceptions, which does not imply $o(X^\delta)$ when $\delta<1$.
